\PassOptionsToPackage{capitalise,nameinlink,noabbrev}{cleveref}
\documentclass[11pt]{article}
\usepackage[margin=1in]{geometry}
\usepackage{amsmath,amssymb}
\usepackage[colorlinks=true,linkcolor=blue,citecolor=blue,urlcolor=blue]{hyperref}
\usepackage{cleveref}
\usepackage{parskip}

\newcommand\blfootnote[1]{%
  \begingroup
  \renewcommand\thefootnote{}\footnote{#1}%
  \addtocounter{footnote}{-1}%
  \endgroup
}

\title{Do Not Vote the Rich Away:\\
Why the Poor Voter Loses by Confiscating the Wealth That Serves Him}

\author{J. Christian McKinley}

\date{May 2026}

\begin{document}

\maketitle
\blfootnote{Zenodo DOI: [to be assigned on deposit].}

\begin{abstract}
This paper addresses the voter who has been told that confiscating the wealth of the rich will improve his condition. The argument is direct: in a property-rights economy, the rich become rich by producing goods and services that others voluntarily purchase, and they remain rich only by continuing to do so. The voter who confiscates is therefore not punishing an enemy; he is shutting down the mechanism by which his own wages, products, and services are funded. The paper does not argue from charity toward the wealthy. It argues from self-interest of the non-wealthy. The structural premise that retention degrades non-holders is denied in \cite{nogo}; this paper draws the political conclusion that follows.
\end{abstract}

\section{Who This Paper Is For}

This paper is written for the voter who is not rich and who has been told that voting away the wealth of the rich will help him. He may have been told this in different vocabularies---fairness, justice, equality, the common good---but the operative claim is the same: there are people with more than they need, and taking from them will leave him with more than he has.

The argument here is that this claim is false on its own terms, by which I mean: false even if one accepts that the voter's only concern is his own welfare. The case against confiscation does not require sympathy for the wealthy. It requires only attention to where the voter's wages, products, and services actually come from.

\section{How the Rich Become Rich}

In a property-rights economy, large fortunes are not seized. They are accumulated, one transaction at a time, from people who chose to part with their money in exchange for something they wanted more than the money.

This is not a moral observation. It is a mechanical one. To grow a business, the owner must persuade strangers to hand him cash voluntarily, repeatedly, in numbers large enough to exceed his costs and his competitors. Every dollar in his account is a dollar that someone else, somewhere, decided was worth less than what he was offering in exchange. If he stops offering value, the dollars stop arriving. If a competitor offers more value, the dollars go there instead.

The fortune is therefore a record of cumulative voluntary judgment by customers. It is not evidence of theft. It is evidence of service.

This is true of the merchant who sells groceries, the manufacturer who builds appliances, the developer who writes software, the entrepreneur who organizes a delivery network, the investor who funds a factory before its product exists. In each case the wealth grows because someone was willing to pay for what was offered, and the willingness arose because what was offered improved the buyer's life by his own assessment.

The voter who is not rich participates in this process every day, on the customer side. He buys what he wants. He passes by what he does not want. The fortunes accumulating around him are the cumulative trace of choices like his, made by millions of people, registering what they preferred to keep their money for.

\section{What About Inheritance?}

The voter is sometimes told that this argument fails for inherited wealth. The heir, after all, did not earn it. He did not serve customers. He simply received what someone else had built. Why should the voter respect a fortune that was not earned by its current holder?

The answer is that the voter's question concerns the present, not the past. Whatever the origin of the wealth, the heir who wishes to preserve it---let alone grow it---is subject to the same discipline as the original earner. There is no fourth option in which inherited wealth sits idle and persists.

Capital decays. Inflation erodes it. Taxation reduces it. Consumption depletes it. Mismanagement loses it to better allocators. The heir who does nothing watches the fortune dissolve. The heir who consumes returns the wealth to the productive economy through his spending, where it funds the activities of others. The heir who wishes the fortune to grow has only one path available: deploy it where it earns a return. Returns, in a property-rights economy, come from customers paying voluntarily for goods and services.

The competent heir therefore serves customers, or hires managers who serve customers, or invests in enterprises that serve customers. The mechanism is identical to the one that produced the wealth in the first place. The lazy heir loses the fortune. The careless heir loses it. The competent heir grows it by continuing the same productive activity the voter benefits from.

The origin of the wealth is therefore irrelevant to the voter's present question. Whether the current holder earned it or inherited it, the wealth persists only by continuing to serve people like him. To confiscate it because the holder did not personally earn it is to interrupt the activity that funds the voter's life on grounds unrelated to whether the activity helps him.

\section{What Confiscation Shuts Down}

When the voter authorizes the seizure of these fortunes---through a wealth tax, a forced redistribution, a confiscatory rate, or any of the softer instruments that perform the same function---he is not removing a parasite from the system. He is removing the signal the system runs on.

The wealth-holder accumulates because he serves customers. If the accumulation is taken, the incentive to serve customers weakens. This is not a moral claim about the wealthy person's character. It is a structural claim about the relationship between effort and reward. When the link is severed, fewer people invest the years required to build a business; fewer people accept the risk required to develop a new product; fewer people maintain operations whose returns will be confiscated before they compound.

The voter does not feel this immediately. The factory does not close the day after the vote. The product does not disappear from the shelf the following week. The compounding loss is invisible at the moment of decision, and visible only years later, in the form of a job that was not created, a product that was not developed, a price that did not fall, a service that did not exist.

By the time the loss is visible, it is no longer attributable. The voter does not see the connection between the policy he supported and the absence of the thing that would otherwise have been built. He sees only that things are harder than they were, and he is told, by the same voices that recommended the original confiscation, that the answer is more of it.

This is the trap. The mechanism that produced the goods and services he relied on was the mechanism he was told to vote against. He cannot tell, after the fact, what would have existed had the mechanism been left alone.

\section{The Confusion the Voter Is Asked to Make}

The political case for confiscation rests on a confusion the voter is encouraged not to examine.

He is told that the wealthy person's fortune is a thing the wealthy person has, separate from him. He is told that taking it will leave the wealthy person with less and the voter with more. He is told that the wealthy person does not need it.

What he is not told is that the fortune exists because the wealthy person did something the voter, or others like him, valued enough to pay for. The fortune is not a stockpile sitting beside the productive economy. It is the productive economy's record of having served customers. To take it is not to redirect a stagnant pool. It is to interrupt the activity that produced the pool, and to discourage the next round of activity that would have produced more of it.

The voter is also not told that the wealthy person, after confiscation, has options. He can move. He can stop. He can invest elsewhere. He can scale back. The voter cannot. The voter depends on the productive activity continuing in his country, in his region, in his town. When the wealthy person responds to confiscation by reducing the activity the voter depends on, it is the voter who absorbs the cost. The wealthy person, if he has resources, has already arranged otherwise.

The voter is told that the wealthy are stuck with him. They are not. He is stuck with the consequences of the system he votes for.

\section{The Reframe}

The standard political question is: \emph{should the rich be allowed to keep what they have?} This phrasing centers the rich person as the subject and the voter as the disinterested judge of whether retention is permitted.

The honest question is different: \emph{is the voter better off when the productive activity around him continues, or when it is interrupted?} This phrasing centers the voter as the subject and the wealth-holder as one component of the activity the voter relies on.

Stated this way, the answer is not in serious doubt. The voter is better off when the supermarket is open, the factory is running, the apartment is being built, the software is being written, the medicine is being developed, the truck is making the delivery. Each of these activities, in a property-rights economy, is undertaken by someone who expects to be compensated and who will reduce or abandon the activity if the compensation is taken.

The voter has been encouraged to focus on the wealth-holder's bank account. The relevant object is the activity the bank account incentivizes. The activity is what the voter needs. The bank account is the mechanism by which the activity is sustained.

\section{Who Profits from the Voter's Confusion}

The argument for confiscation is not delivered to the voter for the voter's benefit. It is delivered by political actors who benefit from the confiscatory machinery itself.

The legislator who votes the confiscation acquires power he did not previously have. The administrator who distributes the seized funds acquires discretion. The party that runs on the platform acquires votes. The intellectual who supplies the moral vocabulary acquires status. None of these actors is the voter, and none of their gains accrue to the voter.

The voter has been recruited into a coalition whose payouts go to its organizers. He is told the cause is justice. The cause is the redistribution of authority from the producer to the political class. Whether any portion of the seized wealth reaches the voter depends on the discretion of actors who answer to political incentives, not to him. Historically, the portion is small, the duration is short, and the long-run cost---in foregone production, foregone employment, foregone goods---is borne by the voter and his children.

The voter who sees this clearly is sometimes accused of defending the rich. He is doing nothing of the sort. He is defending the activity that pays his wages and stocks his shelves, against political actors who profit from disrupting it.

\section{What the Voter Should Ask}

When the voter is presented with a proposal to take wealth from the rich, he should ask three questions.

\emph{First:} What activity, currently being performed by people who expect to be paid, will continue if this passes, and what activity will not? The answer is rarely zero on the second count.

\emph{Second:} Who, specifically, receives the seized funds, and through what mechanism? If the answer is a discretionary fund administered by political actors, the voter should expect that most of the funds will not reach him.

\emph{Third:} What does the proposer gain from passing the policy? If the proposer gains power, votes, status, or career advancement, the voter should weight the proposer's enthusiasm accordingly.

These questions do not require the voter to love the rich. They require only that he ask whether his own life improves under the proposed policy. In most cases, asked seriously, the answer is no.

\section{Conclusion}

The voter who is not rich has been told that voting the rich away will improve his condition. The argument of this paper is that the opposite is true. In a property-rights economy, the wealth of the rich is the record of activity the voter has benefited from, and the continuation of that activity is what the voter depends on. Confiscation interrupts the activity. The voter pays the cost.

The case against confiscation is not made on behalf of the rich. It is made on behalf of the voter who has been encouraged to confuse the destruction of productive activity with his own advancement. The two are not the same. They are opposite.

Do not vote the rich away. They are, by the structure of how they earn---and by the structure of how they hold what they have inherited---working to make your life better. The voice that tells you otherwise is not your ally.

\begin{thebibliography}{9}

\bibitem{nogo}
J. C. McKinley, \emph{Richness Does Not Make the Non-Rich Worse Off: A Structural No-Go}, manuscript (May 2026).

\bibitem{kindness}
J. C. McKinley, \emph{Destructive Kindness: How Redistribution Inverts Incentives and Shrinks Production}, Zenodo (2025). \href{https://doi.org/10.5281/zenodo.17773660}{10.5281/zenodo.17773660}.

\bibitem{legislated}
J. C. McKinley, \emph{Taxation as Legislated Crime via Majority Hegemony}, Zenodo (2026). \href{https://doi.org/10.5281/zenodo.18530711}{10.5281/zenodo.18530711}.

\bibitem{inequality}
J. C. McKinley, \emph{Property Rights and the False Moralization of Inequality}, manuscript (Feb. 2026).

\end{thebibliography}

\end{document}